% problem-type: ham-dieu-hoa-cuc-dai-duy-nhat
% order: 95
% Nguon: De thi MAT2306 (PTDHR 1) cuoi ky 2022-2023 Cau 1 + 2021-2022 Cau 1. (co dap an chinh thuc)

\section{Hàm điều hòa --- cực đại, Liouville \& duy nhất nghiệm}

\subsection*{Đề bài ví dụ}

\paragraph{Ví dụ A (nửa mặt phẳng).}\leavevmode

Cho nửa mặt phẳng trên $H=\{(x,y)\in\mathbb{R}^2 : y>0\}$.

(a) Giả sử hàm liên tục $u:\overline{H}\to\mathbb{R}$ là hàm điều hòa, bị chặn trong $H$ và $u(x,0)=0,\ \forall x\in\mathbb{R}$. Chứng minh rằng $u=0$ trong $\overline{H}$. Nếu bỏ tính bị chặn đi thì đpcm còn đúng không? Tại sao?

(b) Sử dụng kết quả câu (a) để chứng minh tính duy nhất nghiệm bị chặn của bài toán biên Dirichlet đối với phương trình Poisson trong nửa mặt phẳng $H$.

\paragraph{Ví dụ B (miền ngoài hình cầu).}\leavevmode

Với $R>1$, xét miền $\Omega_R=\{(x,y,z)\in\mathbb{R}^3 : 1<x^2+y^2+z^2<R^2\}$ trong không gian. Xác định biên của $\Omega_R$. Phát biểu tường minh nguyên lý cực đại cho hàm $v\in C(\overline{\Omega_R})$ điều hòa trong $\Omega_R$. Từ đó hãy chứng minh bài toán biên Dirichlet ngoài hình cầu sau
\[
\begin{cases}
\Delta u = f & \text{ngoài } \overline{B_1},\\
u = \varphi & \text{trên } \partial B_1,
\end{cases}
\]
trong đó $B_1=\{(x,y,z)\in\mathbb{R}^3 : x^2+y^2+z^2<1\}$, $f\in C(\mathbb{R}^3\setminus\overline{B_1})$, $\varphi\in C(\partial B_1)$ là các hàm cho trước, có tối đa một nghiệm thỏa mãn $\lim_{r\to\infty}u(x,y,z)=0$, với $r=\sqrt{x^2+y^2+z^2}$.

\emph{\small(Nguồn: Đề thi MAT2306 --- PTĐHR 1, cuối kỳ 2022--2023 Câu 1 và 2021--2022 Câu 1. Có đáp án chính thức.)}

\subsection*{Nhận ra dạng này khi nào?}

\begin{itemize}
  \item Đề hỏi chứng minh bài toán \textbf{DUY NHẤT} nghiệm (hoặc tối đa một nghiệm), nhất là nghiệm \emph{bị chặn} hay nghiệm \emph{tiến về $0$ ở vô cùng}.
  \item Đề cho hàm điều hòa \textbf{bị chặn} với dữ kiện biên bằng $0$ và yêu cầu chứng minh hàm $\equiv 0$.
  \item Đề yêu cầu \textbf{phát biểu nguyên lý cực đại} rồi dùng nó trên một miền \emph{không bị chặn} (nửa mặt phẳng, miền ngoài hình cầu).
\end{itemize}

\begin{meobox}
Miền \textbf{không bị chặn} thì không áp dụng nguyên lý cực đại trực tiếp được. Hai mẹo:
\begin{itemize}
  \item Có biên phẳng + dữ kiện biên $0$ $\Rightarrow$ \textbf{thác triển lẻ/chẵn} qua biên rồi dùng định lý Liouville.
  \item Miền ngoài + nghiệm $\to 0$ ở $\infty$ $\Rightarrow$ chặn bằng $\varepsilon$ trên một \textbf{vành} $\Omega_R$, dùng nguyên lý cực đại, rồi cho $R\to\infty$, $\varepsilon\to0$.
\end{itemize}
\end{meobox}

\subsection*{Công thức \& phương pháp}

\begin{nhobox}
\textbf{Nguyên lý cực đại.} Hàm điều hòa trong miền bị chặn, liên tục đến biên thì đạt giá trị lớn nhất và nhỏ nhất \emph{trên biên}. Trên vành $\Omega_R$ (biên gồm hai mặt cầu):
\[
\min\{m_1,m_R\}\le v\le \max\{M_1,M_R\},
\]
với $m_a=\min_{\partial B_a}v$, $M_a=\max_{\partial B_a}v$.

\smallskip
\textbf{Định lý Liouville.} Hàm điều hòa và \emph{bị chặn} trên toàn $\mathbb{R}^n$ thì là hàm hằng.

\smallskip
\textbf{Thác triển qua biên phẳng $y=0$.}
\begin{itemize}
  \item Thác triển \emph{lẻ} (khi $u=0$ trên biên): $U(x,y)=u(x,y)$ nếu $y\ge0$, $=-u(x,-y)$ nếu $y<0$.
  \item Thác triển \emph{chẵn} (khi $u_x=0$ hoặc $u_y=0$ trên biên).
\end{itemize}
Cả hai giữ tính điều hòa qua $y=0$.

\smallskip
\textbf{Duy nhất nghiệm.} Lấy hiệu hai nghiệm $v=u_1-u_2$: $v$ điều hòa, dữ kiện biên $=0$. Chứng minh $v\equiv0$.
\end{nhobox}

\subsection*{Đề bài \& bài giải}

\paragraph{Ví dụ A (nửa mặt phẳng).}\leavevmode

\textbf{(a)} Do $u(x,0)=0$, thác triển \emph{lẻ} hàm $u$ qua trục hoành $y=0$: đặt
\[
U(x,y)=
\begin{cases}
u(x,y) & y\ge 0,\\
-u(x,-y) & y<0.
\end{cases}
\]
Ta được hàm $U$ điều hòa, bị chặn trên toàn mặt phẳng $\mathbb{R}^2$.

Theo định lý Liouville, $U$ là hằng, hay $u$ là hằng. Vì $u(x,0)=0$ nên hằng đó bằng $0$.

\begin{nhobox}
Vậy $u=0$ trên $\overline{H}$.
\end{nhobox}

Nếu không có tính bị chặn thì đpcm \emph{không còn đúng}. Phản ví dụ: lấy $u(x,y)=y$ --- hàm này điều hòa, $u(x,0)=0$ nhưng $u\not\equiv0$.

\textbf{(b)} Giả sử $u_1,u_2$ là hai nghiệm bị chặn của bài toán biên Dirichlet đối với phương trình Poisson trong $H$. Khi đó hiệu $u=u_1-u_2$ là hàm điều hòa (vì hai vế Poisson $\Delta u_1=\Delta u_2=f$ triệt tiêu), bị chặn, và $u=0$ trên $y=0$ --- tức là $u$ thỏa mãn câu (a).

Từ câu (a) ta có $u=0$, hay $u_1=u_2$ trong $\overline{H}$.

\begin{nhobox}
Bài toán Dirichlet cho Poisson trong $H$ có tối đa một nghiệm bị chặn.
\end{nhobox}

\paragraph{Ví dụ B (miền ngoài hình cầu).}\leavevmode

\textbf{Biên của $\Omega_R$.} Miền $\Omega_R=\{1<x^2+y^2+z^2<R^2\}$ có biên gồm hai mặt cầu
\[
\partial B_1=\{x^2+y^2+z^2=1\}
\quad\text{và}\quad
\partial B_R=\{x^2+y^2+z^2=R^2\}.
\]

\textbf{Nguyên lý cực đại.} Cho $v\in C(\overline{\Omega_R})$ điều hòa trong $\Omega_R$. Khi đó
\[
\min\{m_1,m_R\}\le v(x,y,z)\le \max\{M_1,M_R\},\quad \forall(x,y,z)\in\overline{\Omega_R},
\]
trong đó $m_a=\min_{\partial B_a}v$, $M_a=\max_{\partial B_a}v$, $a\in\{1,R\}$.

\textbf{Chứng minh duy nhất.} Giả sử $u_1,u_2$ là các nghiệm của bài toán Dirichlet ngoài hình cầu, đều thỏa $\lim_{r\to\infty}u_j=0$, $j=1,2$. Khi đó hiệu $v=u_1-u_2$ thỏa mãn
\[
\begin{cases}
\Delta v = 0 & \text{ngoài } \overline{B_1},\\
v = 0 & \text{trên } \partial B_1,\\
\lim_{r\to\infty}v(x,y,z)=0. &
\end{cases}
\]

Ta sẽ chứng minh $v(x,y,z)=0$ với mọi $x^2+y^2+z^2>1$, bằng cách chứng minh
\[
|v(x,y,z)|\le \varepsilon,\quad \forall x^2+y^2+z^2>1,\ \forall \varepsilon>0.
\]

Lấy $\varepsilon>0$ và điểm $(x_0,y_0,z_0)\in\mathbb{R}^3\setminus\overline{B_1}$, nghĩa là $x_0^2+y_0^2+z_0^2>1$. Do $\lim_{r\to\infty}v=0$ nên tồn tại $R^2>x_0^2+y_0^2+z_0^2$ sao cho
\[
|v(x,y,z)|\le \varepsilon,\quad \forall x^2+y^2+z^2\ge R^2.
\]
Khi đó trên mặt cầu ngoài: $m_R=\min_{\partial B_R}v\ge -\varepsilon$, $M_R=\max_{\partial B_R}v\le \varepsilon$.

Lại có $v\big|_{\partial B_1}=0$, nên áp dụng nguyên lý cực đại cho hàm điều hòa trên vành $\Omega_R$:
\[
-\varepsilon\le \min\{0,m_R\}\le v(x,y,z)\le \max\{0,M_R\}\le \varepsilon,\quad \forall(x,y,z)\in\overline{\Omega_R}.
\]

Nhắc lại $1<x_0^2+y_0^2+z_0^2<R^2$, nên điểm $(x_0,y_0,z_0)$ nằm trong $\Omega_R$, do đó $|v(x_0,y_0,z_0)|\le \varepsilon$. Cho $\varepsilon\to0$ ta được $v(x_0,y_0,z_0)=0$.

\begin{nhobox}
Vì $(x_0,y_0,z_0)$ tùy ý ngoài $B_1$ nên $v\equiv0$, hay $u_1=u_2$. Bài toán Dirichlet ngoài hình cầu có tối đa một nghiệm thỏa $u\to0$ khi $r\to\infty$.
\end{nhobox}
